\section{Programmkonzept}
\label{sec:programmkonzept}

\subsection{Programmstruktur}
Das Programm besteht aus wenigen Bausteinen, die im \texttt{OB1} nacheinander aufgerufen werden:
\begin{enumerate}[itemsep=0pt, topsep=2pt]
    \item Der FC \texttt{PT100\_conversion} rechnet die Spannung am Analogeingang in eine Temperatur in \textdegree C um und prüft, ob der Messwert gültig ist.
    \item Der FB \texttt{TempDisplay\_Ctrl} mit eigenem Instanz-Datenbaustein enthält die komplette Steuerung. Er bekommt die Temperatur, die Taster und die Taktmerker und setzt die LEDs.
    \item Innerhalb von \texttt{TempDisplay\_Ctrl} wird der vorgegebene \texttt{Motor\_CTRL\_FB} als Multiinstanz aufgerufen. Er steuert den Schrittmotor und wertet den optischen Sensor für die Referenzfahrt aus.
\end{enumerate}
Der rote Taster ist ein Öffner. Sein Signal wird vor dem Aufruf umgedreht, damit am Baustein beide Taster gleich arbeiten (TRUE = gedrückt).

\subsection{Ablaufsteuerung}
Das ganze Programm für die Temperaturanzeige steckt in einem Funktionsbaustein, dem \texttt{TempDisplay\_Ctrl}. 
Er ist als Zustandsautomat aufgebaut: Das Programm befindet sich immer in genau einem Zustand und wechselt bei bestimmten Bedingungen in einen anderen. 
Der Baustein wird in jedem SPS-Zyklus aufgerufen und arbeitet dabei immer dieselben Schritte ab:
Zuerst werden die Taster eingelesen und die Zeiten aktualisiert. Dann wird aus der Temperatur berechnet, wohin der Zeiger zeigen soll. 
Anschließend wird geprüft, ob ein Zustandswechsel nötig ist, der aktuelle Zustand wird bearbeitet, der Motorbaustein \texttt{Motor\_CTRL\_FB} wird aufgerufen und zum Schluss werden die LEDs gesetzt.

Es gibt sechs Zustände:
\begin{itemize}[itemsep=0pt, topsep=2pt]
    \item \textbf{INIT:} Start nach dem Einschalten oder nach einem Neustart. Eine laufende Bewegung wird gestoppt und alle Werte werden zurückgesetzt. Danach geht es direkt in die Referenzfahrt.
    \item \textbf{REF:} Referenzfahrt. Der Zeiger dreht sich, bis der optische Sensor ihn in der Nullposition erkennt.
    \item \textbf{DISPLAY:} Normalbetrieb. Der Motor steht und die Temperatur wird überwacht.
    \item \textbf{MOVE:} Der Zeiger fährt auf den neuen Temperaturwert.
    \item \textbf{JOG:} Der Bediener verstellt den Zeiger von Hand mit den Tastern.
    \item \textbf{ERROR:} Es ist ein Fehler aufgetreten. Der Motor wird sofort gestoppt und die rote LED blinkt.
\end{itemize}

Wohin der Zeiger zeigen soll, ergibt sich aus
\[
\varphi_{\mathrm{Ziel}} = 10 \cdot T + \text{Offset}
\]
Der Wert wird auf \mbox{0 bis 359{,}9\textdegree{}} begrenzt. Bei 25\,\textdegree C zeigt der Zeiger also auf 250. 
Der Offset ist die Korrektur, die der Bediener von Hand eingestellt hat (siehe Kapitel~\ref{sec:korrekturfahrt}). 
Die Zahlen auf der Skala werden gegen den Uhrzeigersinn größer. Ein größerer Wert bedeutet deshalb eine Drehung gegen den Uhrzeigersinn.
Abbildung~\ref{fig:programmablauf} zeigt den Ablauf.

\begin{figure}[H]
\centering
\begin{tikzpicture}[
    every node/.style={font=\footnotesize},
    proc/.style={rectangle, draw, rounded corners=2pt, minimum width=4.6cm, minimum height=0.75cm, align=center, fill=blue!6, inner sep=2pt},
    dec/.style={diamond, draw, aspect=3.4, inner sep=0pt, align=center, fill=orange!10},
    term/.style={rectangle, draw, rounded corners=0.35cm, minimum width=3.4cm, minimum height=0.6cm, align=center, fill=gray!15},
    side/.style={rectangle, draw, rounded corners=2pt, minimum width=3.3cm, minimum height=0.75cm, align=center, inner sep=2pt},
    arr/.style={-{Stealth[length=1.8mm]}, thick},
    lbl/.style={font=\scriptsize, fill=white, inner sep=1pt}
]
% ---------------- Hauptspalte ----------------
\coordinate (R) at (3.5,0);   % rechte Rückweg-Spur
\coordinate (L) at (-3.75,0);  % linke Rückweg-Spur
\node[term] (start) at (0,0) {SPS-Start / Neustart};
\node[proc, below=0.7cm of start] (init)  {\textbf{INIT}: Motor stoppen, Variablen zurücksetzen\\beide LEDs blinken mit 2\,Hz};
\node[proc, below=0.7cm of init]  (ref)   {\textbf{REF}: Referenzfahrt zur Nullposition\\beide LEDs blinken mit 2\,Hz};
\node[dec,  below=0.7cm of ref]   (refok) {Referenz\\gefunden?};
\node[proc, below=0.8cm of refok] (disp)  {\textbf{DISPLAY}: Motor steht, Anzeige aktiv\\grüne LED blinkt mit 1\,Hz};
\node[dec,  below=0.7cm of disp]  (valid) {Temperaturwert\\gültig?};
\node[dec,  below=0.8cm of valid] (jogreq){grün oder rot\\> 300\,ms gedrückt?};
\node[dec,  below=0.8cm of jogreq](upd)   {1-s-Takt oder\\Zwangsupdate?};
\node[dec,  below=0.8cm of upd]   (delta) {$|\Delta \varphi| \ge \text{HYST}$?};
\node[proc, below=0.8cm of delta] (move)  {\textbf{MOVE}: Positionierfahrt auf Zielwinkel\\grüne LED blinkt mit 1\,Hz};
% ---------------- Seitenspalten ----------------
\node[side, fill=blue!6] (jog)   at (jogreq -| -6.0,0) {\textbf{JOG}: Korrekturfahrt\\solange Taster gedrückt\\LED des Tasters leuchtet};
\node[side, fill=red!10] (error) at (valid -| 5.7,0)   {\textbf{ERROR}: Motor gestoppt\\rote LED blinkt mit 1\,Hz};

% ---------------- Hauptablauf ----------------
\draw[arr] (start) -- (init);
\draw[arr] (init) -- (ref);
\draw[arr] (ref) -- (refok);
\draw[arr] (refok) -- node[lbl, right=2pt] {ja} (disp);
\draw[arr] (disp) -- (valid);
\draw[arr] (valid) -- node[lbl, right=2pt] {ja} (jogreq);
\draw[arr] (jogreq) -- node[lbl, right=2pt] {nein} (upd);
\draw[arr] (upd) -- node[lbl, right=2pt] {ja} (delta);
\draw[arr] (delta) -- node[lbl, right=2pt] {ja} (move);

% ---------------- Fehlerpfade (rechts) ----------------
\draw[arr] (refok.east) -| node[lbl, pos=0.35, above] {nein / Timeout 30\,s} (error.north);
\draw[arr] (valid.east) -- node[lbl, above, pos=0.2] {nein} (error.west);
\draw[arr, dashed] (move.east) -| node[lbl, pos=0.45, below] {Motorfehler, Abbruch, Timeout} (error.south);
\draw[arr, dashed] (error.east) -- ++(0.4,0) |- node[lbl, pos=0.75, above] {beide Taster 3\,s} (init.east);

% ---------------- Rückwege zu DISPLAY (rechts) ----------------
\draw[thick] (upd.east) -- (upd.east -| R);
\draw[thick] (delta.east) -- (delta.east -| R);
\draw[arr] (delta.east -| R) -- node[lbl, right, pos=0.3] {nein} ([yshift=-1.5mm]disp.east -| R) -- ([yshift=-1.5mm]disp.east);

% ---------------- Rückwege zu DISPLAY (links) ----------------
\draw[arr] (jogreq.west) -- node[lbl, above, pos=0.15] {ja} (jog.east);
\draw[arr] (jog.north) |- node[lbl, pos=0.75, above] {Taster losgelassen} ([yshift=1.5mm]disp.west);
\draw[arr] (move.west) -- (move.west -| L) -- node[lbl, left, pos=0.08] {Position erreicht} ([yshift=-1.5mm]disp.west -| L) -- ([yshift=-1.5mm]disp.west);
\end{tikzpicture}
\caption{Programmablauf des Zustandsautomaten \texttt{TempDisplay\_Ctrl}}
\label{fig:programmablauf}
\end{figure}

Zwei Dinge gelten in jedem Zustand:
Werden beide Taster 3\,s gleichzeitig gehalten, startet das Programm neu. Fehler und Offset werden gelöscht, und es geht wieder bei \texttt{INIT} los, also mit einer neuen Referenzfahrt. 
Meldet der Motorbaustein einen Fehler, wechselt das Programm sofort in den Zustand \texttt{ERROR}.
Außerdem prüft das Programm im Normalbetrieb, ob die Nullposition noch bekannt ist. Ist sie verloren gegangen, wird automatisch eine neue Referenzfahrt gestartet.
Bei jedem Fehler bleibt der Motor stehen und die rote LED blinkt.

\subsection{Positionierfahrt}
Im Normalbetrieb wird der Zeiger einmal pro Sekunde an die aktuelle Temperatur angepasst. Direkt nach einer Referenzfahrt passiert das sofort, ohne auf die nächste Sekunde zu warten. 
Dazu berechnet das Programm, wie weit der Zeiger vom richtigen Wert entfernt ist:
\[
\Delta \varphi = \varphi_{\mathrm{Ziel}} - \varphi_{\mathrm{Ist}}
\]

Der Zeiger fährt nur los, wenn diese Abweichung mindestens so groß ist wie ein fest eingestellter Mindestwert (Hysterese \texttt{HYST\_ANGLE}). 
So fährt der Motor nicht bei jeder winzigen Schwankung der Temperatur hin und her.

Für die Fahrt bekommt der Motorbaustein den Winkel, um den er sich drehen soll. Die Drehrichtung ergibt sich daraus, ob die Temperatur gestiegen oder gefallen ist. 
Dann wird der Motor gestartet, und das Programm wartet im Zustand \texttt{MOVE}, bis der Motor meldet, dass er angekommen ist. Danach geht es zurück in den Normalbetrieb. 
Dauert die Fahrt zu lange oder wird sie abgebrochen, wechselt das Programm in den Fehlerzustand (Abbildung~\ref{fig:positionierfahrt}).

\begin{figure}[H]
\centering
\begin{tikzpicture}[
    node distance=0.58cm,
    every node/.style={font=\small},
    proc/.style={rectangle, draw, rounded corners=2pt, minimum width=5.5cm, minimum height=0.8cm, align=center, fill=blue!6},
    dec/.style={diamond, draw, aspect=2.5, inner sep=1.2pt, align=center, fill=orange!10},
    term/.style={rectangle, draw, rounded corners=0.35cm, minimum width=3.9cm, minimum height=0.75cm, align=center, fill=gray!15},
    arr/.style={-{Stealth[length=2mm]}, thick},
    lbl/.style={font=\footnotesize}
]
\node[term] (s) {DISPLAY: 1-s-Takt\\oder Zwangsupdate};
\node[proc, below=of s] (d) {Abweichung berechnen\\$\Delta \varphi = \varphi_{\text{Ziel}} - \varphi_{\text{Ist}}$};
\node[dec, below=of d] (h) {$|\Delta \varphi| \ge \text{HYST\_ANGLE}$?};
\node[proc, below=0.75cm of h] (p) {diskreter Modus\\Winkel $|\Delta \varphi|$, Richtung aus Vorzeichen};
\node[proc, below=of p] (st) {\textbf{MOVE}: Startimpuls an \texttt{Motor\_CTRL\_FB}};
\node[dec, below=of st] (run) {Motor läuft\\(\texttt{done = 0})?};
\node[dec, below=0.75cm of run] (done) {Position erreicht\\(\texttt{done = 1})?};
\node[term, below=0.75cm of done] (e) {zurück zu DISPLAY};
\node[rectangle, draw, rounded corners=2pt, minimum width=2.6cm, minimum height=0.75cm, align=center, fill=red!10, left=1.4cm of done] (err) {\textbf{ERROR}\\rote LED blinkt};

\draw[arr] (s) -- (d);
\draw[arr] (d) -- (h);
\draw[arr] (h) -- node[lbl, right] {ja} (p);
\draw[arr] (p) -- (st);
\draw[arr] (st) -- (run);
\draw[arr] (run) -- node[lbl, right] {ja} (done);
\draw[arr] (done) -- node[lbl, right] {ja} (e);
\draw[arr] (done.west) -- node[lbl, above] {Abbruch} (err.east);

\draw[arr] (h.east) -- ++(2.2,0) |- node[lbl, right, pos=0.25] {nein} (e.east);
\draw[arr] (run.east) -- ++(0.9,0) |- node[lbl, right, pos=0.25] {nein} ($(st.south)!0.5!(run.north)$);
\draw[arr] (done.east) -- ++(0.9,0) |- node[lbl, right, pos=0.25] {nein} ($(run.south)!0.5!(done.north)$);
\end{tikzpicture}
\caption{Ablauf einer automatischen Positionierfahrt}
\label{fig:positionierfahrt}
\end{figure}

\subsection{Korrekturfahrt}
\label{sec:korrekturfahrt}
Mit der Korrekturfahrt kann der Bediener den Zeiger von Hand verstellen, zum Beispiel wenn die Anzeige nicht ganz stimmt.
Dazu muss er den grünen oder roten Taster länger als 300\,ms gedrückt halten. Ein kurzes oder versehentliches Antippen löst also nichts aus.

Solange der grüne Taster gehalten wird, fährt der Zeiger zu größeren Werten. Solange der rote Taster gehalten wird, fährt er zu kleineren Werten. 
Am Anfang und am Ende der Skala bleibt der Zeiger automatisch stehen.

Lässt der Bediener den Taster los, wartet das Programm, bis der Motor steht. Dann wird gespeichert, wie weit der Zeiger jetzt vom berechneten Wert abweicht:
\[
\text{Offset} = \varphi_{\mathrm{Ist}} - 10 \cdot T
\]
Dieser Offset ist nach oben und unten begrenzt (\texttt{OFFSET\_LIMIT}) und wird nur gespeichert, wenn die Temperatur im Anzeigebereich liegt. 
Ab jetzt wird der Offset bei jeder Anzeige mitgerechnet. Der Zeiger wird danach sofort neu positioniert.

Während der Korrektur leuchtet die LED des Tasters, der gerade gedrückt ist.

% Belegung der Taster und Bedeutung der LEDs
% Wird aus 05_Programmkonzept.tex eingebunden
\subsection{Belegung der Taster}
\label{sec:bedienkonzept}
Zur Bedienung gibt es nur die zwei Taster des AC2398. 
Damit kann der Bediener den Zeiger korrigieren und das Programm nach einem Fehler neu starten. 
Der rote Taster ist ein Öffner, sein Signal wird deshalb vor dem Baustein umgedreht.

\begin{table}[H]
\centering\small
\begin{tabular}{|L{3.6cm}|L{10.4cm}|}
\hline
\textbf{Bedienung} & \textbf{Funktion} \\ \hline
Grün länger als 300\,ms halten & Zeiger fährt nach oben (positiver Offset), bis der Taster losgelassen oder das Skalenende erreicht wird. Grüne LED leuchtet. \\ \hline
Rot länger als 300\,ms halten & Zeiger fährt nach unten (negativer Offset), bis der Taster losgelassen oder das Skalenende erreicht wird. Rote LED leuchtet. \\ \hline
Beide Taster 3\,s halten & Neustart mit Referenzfahrt: Offset wird gelöscht, der Zeiger fährt auf 0 und anschließend auf die aktuelle Temperatur. Beide LEDs blinken mit 2\,Hz, solange die Referenzfahrt läuft. Ein anliegender Fehler wird dabei quittiert. \\ \hline
\end{tabular}
\caption{Belegung der Taster des AC2398}
\end{table}

\subsection{Bedeutung der LEDs}
Die LEDs zeigen über schnelles oder langsames Blinken an, was das Programm gerade macht. 
Wird ein Taster gedrückt, leuchtet seine LED dauerhaft. Das gilt nicht während des Starts, der Referenzfahrt und im Fehlerzustand.

\begin{table}[H]
\centering\small
\begin{tabular}{|L{2.6cm}|L{2.6cm}|L{8.8cm}|}
\hline
\textbf{LED grün} & \textbf{LED rot} & \textbf{Bedeutung} \\ \hline
blinkt 2\,Hz & blinkt 2\,Hz & Referenzfahrt läuft \\ \hline
blinkt 1\,Hz & aus & Normalbetrieb, Anzeige aktuell \\ \hline
blinkt 1\,Hz & blinkt 1\,Hz & Temperatur außerhalb 0--36\,\textdegree C, Zeiger am Skalenende \\ \hline
aus & blinkt 1\,Hz & Fehler, Motor gestoppt \\ \hline
\end{tabular}
\caption{Bedeutung der LED-Zustände}
\end{table}

Eine kompakte Anleitung für den Bediener findet sich in Kapitel~\ref{sec:bedienungsanleitung}.


% Signalliste
\section{Signalliste}
% Tabelle: Name, Datentyp, Ort (PLC-Tag/E/A, Merker, DB), Adresse, Beschreibung
% getrennt nach Eingängen, Ausgängen, internen Variablen und Schnittstellen zum \texttt{Motor\_CTRL\_FB}

